%======================================================================
% Appendix A -- Foundations. Written to be \input by paper.tex (its \section
% header and intro live there). These are the derivations the NEW results of
% Appendix B take as given; each subsection ends with a pointer to where it is
% used. The source paper's variance decomposition and OOS metrics are recorded
% at the end for reference (they are not prerequisites for the new results).
%======================================================================

\subsection{Reduced form and companion form}\label{app:redcomp}
The starting point is the reduced-form VAR
\begin{equation}
y_t=\nu+A_1y_{t-1}+A_2y_{t-2}+\cdots+A_py_{t-p}+u_t,
\qquad \E[u_t]=0,\ \ \E[u_tu_t\tp]=\Sigma_u,
\tag{1}\label{aeq:red}
\end{equation}
a $K$-dimensional linear autoregression of order $p$. ``Reduced form'' means
$u_t$ collects statistical innovations with \emph{no} economic interpretation;
the $A_i$ summarise own- and cross-dynamics. Because every equation shares the
same right-hand-side regressors, the seemingly-unrelated-regression system
collapses to equation-by-equation OLS.

\paragraph{Companion (state-space) form.} Stacking the current and $p-1$ most
recent vectors turns \eqref{aeq:red} into a first-order system,
\begin{equation}
Y_t=\mu+AY_{t-1}+U_t,
\tag{2}\label{aeq:comp}
\end{equation}
with
\[
Y_t=\begin{bmatrix} y_t\\ y_{t-1}\\ \vdots\\ y_{t-p+1}\end{bmatrix},\quad
\mu=\begin{bmatrix}\nu\\0\\\vdots\\0\end{bmatrix},\quad
U_t=\begin{bmatrix}u_t\\0\\\vdots\\0\end{bmatrix},\quad
A=\begin{bmatrix}
A_1 & A_2 & \cdots & A_{p-1} & A_p\\
I_K & 0 & \cdots & 0 & 0\\
0 & I_K & \cdots & 0 & 0\\
\vdots & & \ddots & & \vdots\\
0 & 0 & \cdots & I_K & 0
\end{bmatrix}.
\]
Reading \eqref{aeq:comp} block-row by block-row: the first block row is exactly
\eqref{aeq:red}; block row $i\ge2$ is the identity $y_{t-i+1}=y_{t-i+1}$, which
merely carries lags forward in the state. With the top-block selector
$J=[\,I_K\ 0\ \cdots\ 0\,]\in\R^{K\times Kp}$ one has the identities used
repeatedly below,
\begin{equation}
y_t=JY_t,\qquad \mu=J\tp\nu,\qquad U_t=J\tp u_t.
\tag{$\star$}\label{aeq:star}
\end{equation}
Covariance-stationarity requires all eigenvalues of $A$ to lie inside the unit
circle.

\medskip
\noindent\textit{Used in:} everything below; the companion matrix $A$ is exported
to the dashboard so the counterfactual baseline (\S\ref{app:cf}) can be iterated
in the browser from any in-sample origin.

\subsection{Forward iteration: the \texorpdfstring{$h$}{h}-step recursion}\label{app:recursion}
Iterating \eqref{aeq:comp} forward $h$ periods gives
\begin{equation}
Y_{t+h}=\Bigg(\sum_{j=0}^{h-1}A^{j}\Bigg)\mu+A^{h}Y_t+\sum_{j=0}^{h-1}A^{j}U_{t+h-j}.
\tag{3}\label{aeq:recursion}
\end{equation}

\paragraph{Derivation (induction on $h$).} The base case $h=1$ is
\eqref{aeq:comp}. Assuming \eqref{aeq:recursion} at horizon $h$ and applying
\eqref{aeq:comp} once more, $Y_{t+h+1}=\mu+AY_{t+h}+U_{t+h+1}$, substitution gives
\[
Y_{t+h+1}
=\underbrace{\mu+\sum_{j=0}^{h-1}A^{j+1}\mu}_{=\sum_{j=0}^{h}A^{j}\mu}
+A^{h+1}Y_t
+\underbrace{\sum_{j=0}^{h-1}A^{j+1}U_{t+h-j}+U_{t+h+1}}_{=\sum_{j=0}^{h}A^{j}U_{t+h+1-j}},
\]
where the last bracket re-indexes $j\mapsto j+1$ (turning the sum into
$\sum_{j=1}^{h}A^{j}U_{t+h+1-j}$) and the stray $U_{t+h+1}$ supplies the missing
$j=0$ term. This is \eqref{aeq:recursion} at $h+1$.

\medskip
\noindent\textit{Used in:} the baseline (\S\ref{app:base}) is its conditional
mean; the moving-average map (\S\ref{app:ma}) is its stochastic part.

\subsection{The baseline forecast}\label{app:base}
Taking $\E[\,\cdot\mid\Fset_t]$ of \eqref{aeq:recursion} annihilates every future
shock, since $U_{t+h-j}$ is dated $t+h-j\ge t+1$ for $j=0,\dots,h-1$ and hence
unforecastable, $\E[U_{t+h-j}\mid\Fset_t]=0$. Premultiplying by $J$ and using
$\mu=J\tp\nu$ from \eqref{aeq:star},
\begin{equation}
\E[y_{t+h}\mid\Fset_t]
=\Bigg(\sum_{j=0}^{h-1}JA^{j}J\tp\Bigg)\nu+JA^{h}Y_t.
\tag{4}\label{aeq:base}
\end{equation}
Under quadratic loss this conditional expectation is the optimal point forecast;
it is the \emph{baseline scenario} against which every conditional forecast and
every counterfactual is measured.

\medskip
\noindent\textit{Used in:} the constraint right-hand side
$r=(\text{target}-\text{baseline})$ of the main text~\eqref{eq:restriction} and the counterfactual
origin baseline $F_{\text{base}}$ (\S\ref{app:cf}).

\subsection{Forecast error, impulse responses, and the stacked map}\label{app:ma}
Subtracting \eqref{aeq:base} from \eqref{aeq:recursion} leaves only the shock
term; projecting onto the top block with $U_{t+h-j}=J\tp u_{t+h-j}$ gives the
deviation of any realised path from the baseline,
\begin{equation}
y_{t+h}-\E[y_{t+h}\mid\Fset_t]
=\sum_{j=0}^{h-1}JA^{j}J\tp\,u_{t+h-j}
=\sum_{j=0}^{h-1}\Phi_j\,u_{t+h-j},
\qquad \Phi_j:=JA^{j}J\tp.
\tag{5}\label{aeq:dev}
\end{equation}
The $\Phi_j$ are the reduced-form moving-average (impulse-response) coefficients,
with $\Phi_0=I_K$. \textbf{Economic reading:} forcing a variable onto a scenario
path (left-hand side $\neq0$) is \emph{equivalent} to imposing a particular
realisation of future innovations (right-hand side); the scenario cannot happen
``for free,'' and the same innovations move every other variable through $\Phi_j$.

\paragraph{The stacked map.} Collect \eqref{aeq:dev} for $h=1,\dots,H$. Writing
$\Delta y=[\,(y_{t+1}-\E)\tp,\dots,(y_{t+H}-\E)\tp\,]\tp\in\R^{KH}$ and stacking
the innovations conformably, the deviation is a single linear image of the future
shocks,
\begin{equation}
\Delta y=\Mmap_u\,u,\qquad
\Mmap_u=\begin{bmatrix}
\Phi_0 & & &\\
\Phi_1 & \Phi_0 & &\\
\vdots & & \ddots &\\
\Phi_{H-1} & \Phi_{H-2} & \cdots & \Phi_0
\end{bmatrix},
\tag{5$'$}\label{aeq:stack}
\end{equation}
a block-lower-triangular $KH\times KH$ matrix (block $(h,s)=\Phi_{h-s}$ for
$s\le h$, zero otherwise). Once a structural impact matrix $D$ is identified
(\S\ref{app:ident2}), the same construction with $\Theta_j:=\Phi_jD$ in place of
$\Phi_j$ gives the \emph{structural} stacked map $\Mmap$ of the main text
\eqref{eq:Mmap}, whose diagonal block is $\Theta_0=D$.

\medskip
\noindent\textit{Used in:} $\Mmap$ is \emph{the} object of the paper. The
restriction rows $R$ are selected rows of $\Mmap$; the propagation
$\Delta y=\Mmap\beps$ delivers every conditional forecast, band
(\S\ref{app:shock}, \S\ref{app:param}), and counterfactual (\S\ref{app:cf}). The
coefficients $\Theta_j$ are exported to the dashboard, which rebuilds $\Mmap$ in
the browser.

\subsection{The conditioning restriction and the minimum-norm shock path}\label{app:minnorm}
A scenario fixes target values for a subset of variables over some horizons.
Reading the corresponding rows off \eqref{aeq:stack} yields a linear system in the
stacked future innovations,
\begin{equation}
R\,u=r,
\tag{5$''$}\label{aeq:Rur}
\end{equation}
where $r\in\R^{m}$ is the gap between the scenario values and the baseline
forecasts on the $m$ constrained cells and $R\in\R^{m\times KH}$ collects the
associated rows of the MA coefficients. Because a scenario fixes far fewer cells
than there are free shocks ($m<KH$), \eqref{aeq:Rur} is underdetermined; WZ, after
Doan--Litterman--Sims, select the minimum-norm solution.

\paragraph{Derivation (constrained minimisation).} Solve
$\min_u\tfrac12u\tp u$ s.t.\ $Ru=r$. The Lagrangian
$\mathcal L=\tfrac12u\tp u+\lambda\tp(r-Ru)$ has first-order conditions
$u=R\tp\lambda$ and $Ru=r$; substituting, $RR\tp\lambda=r$, so
$\lambda=(RR\tp)^{-1}r$ (full row rank $m$) and
\begin{equation}
\hat u=R\tp(RR\tp)^{-1}r=\pinv{R}\,r.
\tag{5$'''$}\label{aeq:minnorm}
\end{equation}
Adding $\hat u$ back through \eqref{aeq:stack} delivers the conditional forecast
for every variable.

\begin{remark}[dimensions of the inverse]
With $R$ of size $m\times KH$ and $m<KH$, $RR\tp$ ($m\times m$) is invertible
while $R\tp R$ ($KH\times KH$) is singular. A source-paper display that writes
$\hat u=R\tp(R\tp R)^{-1}r$ should read $(RR\tp)^{-1}$ as in \eqref{aeq:minnorm};
the two coincide only under the transposed storage convention for $R$. The
substance is unchanged.
\end{remark}

\begin{remark}[``most likely'' shocks]
If $u\sim N(0,I_{KH})$ then \eqref{aeq:minnorm} is the mode (equivalently the mean)
of $u\mid Ru=r$: minimising $u\tp u$ maximises the Gaussian density. This is why
$\hat u$ is the most likely combination of shocks consistent with the scenario ---
the reading the uncertainty derivations of Appendix~\ref{app:new} make precise.
\end{remark}

\medskip
\noindent\textit{Used in:} \eqref{aeq:minnorm} is the object every new derivation
extends --- soft conditioning replaces $(RR\tp)^{-1}$ by $(RR\tp+\Om)^{-1}$
(\S\ref{app:soft}); the shock band conditions the Gaussian around $\hat\beps$
(\S\ref{app:shock}); the counterfactual perturbs a realised $\hat u$
(\S\ref{app:cf}). The structural-space version $\hat\beps=\pinv{R}r$ used in the
main text is \eqref{aeq:minnorm} with $u=D\eps$ (\S\ref{app:norm}).

\subsection{Structural identification}\label{app:ident2}
The reduced-form innovation is decomposed as $u_t=D\eps_t$ with orthonormal
structural shocks,
\begin{equation}
\E[\eps_t]=0,\quad \E[\eps_t\eps_t\tp]=I_K
\ \Longrightarrow\ \Sigma_u=DD\tp.
\tag{6}\label{aeq:struct}
\end{equation}

\paragraph{Order condition and Cholesky.} A symmetric $\Sigma_u$ has $K(K+1)/2$
distinct entries while $D$ has $K^2$; the equations $\Sigma_u=DD\tp$ fall short by
$K^2-K(K+1)/2=K(K-1)/2$ restrictions. Requiring $D$ lower-triangular zeros out
exactly its $K(K-1)/2$ upper entries, closing the gap, and with a positive
diagonal the solution is the unique Cholesky factor
$D=\chol(\hat\Sigma_u)$, with $\hat\Sigma_u=\sum_t\hat u_t\hat u_t\tp/(T-Kp-1)$. A
lower-triangular $D$ means shock $k$ moves variables $1,\dots,k-1$ only with a
lag but $k,\dots,K$ contemporaneously, so the recursive ordering of $y_t$ is
material (global variables first, sticky prices ahead of the activity/policy
block, forward-looking financial variables last). The structural impulse
responses are $\Theta_j=\Phi_jD$.

\medskip
\noindent\textit{Used in:} $D=\Theta_0$ is the diagonal block of $\Mmap$ and the
change of variables $u=D\eps$ that turns the WZ norm into the Mahalanobis norm
(\S\ref{app:norm}); the whole apparatus is identification-agnostic --- replace
the Cholesky $D$ by a sign-restricted or instrument-based $D$ and every formula in
Appendix~\ref{app:new} is unchanged.

\subsection{Reference: variance decomposition and out-of-sample metrics}\label{app:refresults}
The following source-paper results are recorded for completeness; they are
\emph{not} prerequisites for the new derivations of Appendix~\ref{app:new}.

\paragraph{Forecast-error variance decomposition.} With the structural forecast
error $y_{t+h}-\E y_{t+h}=\sum_{j=0}^{h-1}\Theta_j\eps_{t+h-j}$ and orthonormal
$\eps$, the mean squared prediction error and its per-shock split are
\begin{equation}
\MSPE(h)=\sum_{j=0}^{h-1}\Theta_j\Theta_j\tp,\qquad
\MSPE_{k,l}(h)=\sum_{j=0}^{h-1}\theta_{k,l,j}^2,
\tag{7$'$,\,7$''$}\label{aeq:fevd}
\end{equation}
where $\theta_{k,l,j}=[\Theta_j]_{kl}$; the reported share is
$\MSPE_{k,l}(h)/\sum_{l}\MSPE_{k,l}(h)$, and orthogonality is what makes the shares
sum to one across $l$.

\paragraph{Out-of-sample evaluation (planned exercise).} Point accuracy uses an
AR($p$) benchmark --- \eqref{aeq:red} with diagonal $A_i$ --- and $h$-step errors
$e^{(h,m)}_{k,t}$ aggregated into
\begin{equation}
\RMSE(k,h,m)=\sqrt{\tfrac{1}{\bar t-\underline t}\textstyle\sum_t\big(e^{(h,m)}_{k,t}\big)^2},
\qquad
\MAE(k,h,m)=\tfrac{1}{\bar t-\underline t}\textstyle\sum_t\big|e^{(h,m)}_{k,t}\big|,
\tag{9,\,10}\label{aeq:rmse}
\end{equation}
each reported relative to the benchmark (a ratio below one favours the VAR).
Density accuracy uses the continuous ranked probability score in its quantile
(``pinball'') form: with $\rho_\tau(u)=u(\tau-\Id(u<0))$ and predictive quantiles
$\hat q^{(h,m)}_{\tau,k,t}$ on a grid $\tau_1<\cdots<\tau_N$,
\begin{equation}
\CRPS(k,h,m)=\frac{2}{N-1}\sum_{n=1}^{N}
\Bigg(\frac{1}{\bar t-\underline t}\sum_{t}\rho_{\tau_n}\!\big(y_{k,t}-\hat q^{(h,m)}_{\tau_n,k,t}\big)\Bigg),
\tag{11}\label{aeq:crps}
\end{equation}
the trapezoidal discretisation of
$\CRPS(F,y)=2\int_0^1\rho_\tau(y-F^{-1}(\tau))\,d\tau$, which needs only the
predictive quantiles, not the full density. These metrics drive the planned
backtesting harness (\S\ref{sec:oos}).

\subsection*{Summary of the foundation chain}
\[
\underbrace{(1)\Rightarrow(2)}_{\text{companion}}
\ \Rightarrow\
\underbrace{(3)}_{\text{iterate}}
\ \Rightarrow\
\begin{cases}
\E[\cdot\mid\Fset_t]\ \Rightarrow\ (4)\ \text{baseline}\\[2pt]
\text{deviation}\ \Rightarrow\ (5),(5')\ \Rightarrow\ \Mmap\ \text{(stacked map)}\\[2pt]
Ru=r\ \Rightarrow\ \hat u=\pinv{R}r\ \text{(min-norm, WZ)}\\[2pt]
u=D\eps,\ \Sigma_u=DD\tp\ \Rightarrow\ \Theta_j=\Phi_jD\ \text{(identification)}
\end{cases}
\]
Appendix~\ref{app:new} builds on these objects: it adds a distribution around
$\hat u$ (bands), a noisy version of $Ru=r$ (soft conditions), extra rows (joint /
partial-horizon scenarios), and a perturbation of a realised $\hat u$
(counterfactuals).
